\section{Explicit homology functionals over a finite field}
\label{sf-section}

We construct the finite-field functional used to evaluate
the arithmetic cycles.  The construction uses a cyclic
two-cocycle and a sum over projective lines.  In particular,
this step does not require a theorem about the algebraic
\(K\)-groups of finite fields.

We use the bar chains and boundary of
Section~\ref{sh-section}.  For a commutative coefficient
ring \(A\), a degree-\(m\) cochain is an \(A\)-linear
map \(C_m(G;A)\to A\); it is a cocycle if it vanishes
on boundaries.  The group acts trivially on \(A\).

\subsection{A cyclic cocycle with an exact value}

Let \(k\) be a finite field of cardinality \(Q\),
let \(q\ge2\), and let \(\zeta\in k^\times\) have
order \(q\).  Put \(R=\mathbb Z/q\mathbb Z\), and
let \(v:k^\times\to(R,+)\) be a homomorphism.
For \(x\in\mathbb Z\), write
\(\chi_x:C_q\to k^\times\) for
\(\chi_x(j)=\zeta^{xj}\), where \(C_q=\mathbb Z/q\mathbb Z\).
For \(E\ge0\), define
\[
 z_E=\sum_{j_1,\ldots,j_E\in C_q}
             [j_1|1|\cdots|j_E|1]\in C_{2E}(C_q;R).
\]
When \(E=0\), this means the single empty tuple.
For \(E\ge1\), it is the cycle proved in
Lemma~\ref{sh-cycle}.

\begin{lemma}\label{sf-carry}
There is a symmetric integer-valued function
\(u_0:k^\times\times k^\times\to\mathbb Z\) such that
\[
 u_0(b,c)-u_0(ab,c)+u_0(a,bc)-u_0(a,b)=0
\]
and, writing \(u\) for its reduction modulo \(q\),
\[
                \sum_{j\in C_q}u(\zeta^{xj},\zeta^x)=x
                         \qquad\text{in }R.
\]
\end{lemma}
\begin{proof}
The multiplicative group of a finite field is cyclic.
Here is the elementary reason: if \(m\) is the least
common multiple of its element orders, every element is
a root of \(X^m-1\), so its cardinality is at most \(m\).
Lagrange's theorem gives the reverse inequality.
In a finite abelian group, an element of order \(m\)
is obtained by taking, for each prime, a component of
largest prime-power order and multiplying these components.
The component orders are relatively prime.
Thus there is a generator of \(k^\times\).

Put \(N=Q-1=qr\).  We may choose a generator \(g\)
with \(g^r=\zeta\).
Indeed, for an initial generator \(h\), write
\(\zeta=h^{ra}\), where \((a,q)=1\).
Chinese remainders give an integer \(b\equiv a\pmod q\)
which is \(1\) modulo each prime dividing \(N\) but
not \(q\).  Then \((b,N)=1\), and \(g=h^b\) works.

For \(s\in k^\times\), let \(a(s)\in\{0,\ldots,N-1\}\)
be its exponent relative to \(g\), and define
\[
                  u_0(s,t)
                       =\frac{a(s)+a(t)-a(st)}{N}.
\]
The numerator is divisible by \(N\).
Symmetry follows from commutativity, and the displayed
cocycle identity follows by substituting this formula:
every exponent cancels.

Set \(s=\zeta^x\).  Summing over \(j\in C_q\),
the terms \(a(s^j)\) and \(a(s^{j+1})\) cancel,
even when \(s\) has order less than \(q\).  Hence
\[
       \sum_{j\in C_q}u_0(s^j,s)
             =\frac{q\,a(s)}{N}=\frac{a(s)}r.
\]
Since \(g^r=\zeta\), the last integer is the least
nonnegative residue of \(x\) modulo \(q\).
Reducing modulo \(q\) proves the assertion.
\end{proof}

For commuting entries \(h_1,\ldots,h_m\) and an element
\(g\) commuting with them, define the insertion chain
\[
 P_g[h_1|\cdots|h_m]
   =\sum_{i=0}^m(-1)^i
       [h_1|\cdots|h_i|g|h_{i+1}|\cdots|h_m].
\]
This is the source's prism on such chains; its conjugations
disappear because of the stated commutation.

\begin{proposition}\label{sf-cyclic-functional}
For every \(E\ge0\), there is a cocycle
\(t:C_{2E+1}(k^\times;R)\to R\) satisfying
\[
              t(P_g(\chi_x)_*z_E)=x^E v(g)
            \qquad(x\in\mathbb Z,\ g\in k^\times).
\]
\end{proposition}
\begin{proof}
Using Lemma~\ref{sf-carry}, set
\[
 t([h_1|\cdots|h_{2E+1}])
       =\prod_{a=1}^E u(h_{2a-1},h_{2a})\,
                                      v(h_{2E+1}).
\]
This is the cup product \(u^E\smile v\).
For completeness, if \(f\) has degree \(a\) and
\(h\) has degree \(b\), their cup product on a tuple
is \(f\) of its first \(a\) entries times \(h\) of
its last \(b\) entries.
The bar differential gives
\[
       \delta(f\smile h)
          =(\delta f)\smile h+(-1)^a f\smile\delta h.
\]
To check this identity, expand both sides by the faces:
the faces before the dividing position are those of
\(\delta f\), and those after it are those of
\(\delta h\), shifted by \(a\) positions.
The last face of the first group and the first face
of the second group cancel with opposite signs;
the remaining common face multiplies the two entries
at the dividing position.  These are exactly all
the faces of the tuple on the left.
The cocycle identity of Lemma~\ref{sf-carry} and
\(v(ab)=v(a)+v(b)\) therefore imply \(\delta t=0\).

Evaluate \(t\) on an insertion chain of even length.
The terms inserting \(g\) just before and just after
\(h_{2a-1}\) cancel, because \(u\) is symmetric.
Only insertion in the last position remains.  Thus
\[
 t(P_g[h_1|\cdots|h_{2E}])
       =\prod_{a=1}^E u(h_{2a-1},h_{2a})\,v(g).
\]
For \((\chi_x)_*z_E\), the \(E\) summation indices
are independent, so the value is
\[
       \left(\sum_{j\in C_q}u(\zeta^{xj},\zeta^x)\right)^E
                  v(g)=x^Ev(g).
\]
This includes the empty product \(E=0\).
\end{proof}

\subsection{Transfer over projective lines}

We give the chain construction, including the independence
of its choices.  Let a group \(G\) act on a finite set
\(X\) on the right, and let \(M\) be an abelian group.
Suppose
\[
 c(gh,x)=c(g,x)c(h,xg),\qquad c(1,x)=1,
\]
for a function \(c:G\times X\to M\).
For a tuple \([g_1|\cdots|g_m]\), put
\(x_i=xg_1\cdots g_i\), \(x_0=x\), and define
\[
 T_c[g_1|\cdots|g_m]
    =\sum_{x\in X}
       [c(g_1,x_0)|c(g_2,x_1)|\cdots|c(g_m,x_{m-1})].
\]
Extend this map linearly over any commutative ring \(A\).

\begin{lemma}\label{sf-transfer}
The maps \(T_c\) commute with the bar boundary.
If
\[
                c'(g,x)=b(x)^{-1}c(g,x)b(xg)
\]
for a function \(b:X\to M\), then \(T_{c'}\) and
\(T_c\) are chain homotopic.  In particular they give
the same value on cycles after applying a cocycle.
\end{lemma}
\begin{proof}
For an internal face, the two adjacent entries in
the image multiply to \(c(g_ig_{i+1},x_{i-1})\),
by the cocycle identity.  The last face deletes the
same entry on either side.  For the first face,
replace the summation index \(x\) by \(xg_1\).
This is a permutation of \(X\), so the first faces
also agree.  This proves the chain-map assertion.

Here is an explicit homotopy.  On a degree-\(m\) tuple,
write \(c_i=c(g_i,x_{i-1})\),
\(c'_i=c'(g_i,x_{i-1})\), and \(b_i=b(x_i)\).
Set
\[
 K_m[g_1|\cdots|g_m]
 =\sum_{x\in X}\sum_{i=0}^m(-1)^i
          [c_1|\cdots|c_i|b_i|c'_{i+1}|\cdots|c'_m].
\]
Then
\[
                    dK_m+K_{m-1}d=T_{c'}-T_c.
\]
Indeed the faces not adjacent to the inserted entry
are precisely the corresponding terms of
\(-K_{m-1}d\), after the first-face permutation of
\(X\) already described.  The two faces adjacent
to successive inserted entries cancel by
\(c_i b_i=b_{i-1}c'_i\).
The two unpaired end faces are the all-\(c'\) tuple
and minus the all-\(c\) tuple.
This accounts for every face and its sign.
For a cycle \(z\), the identity expresses
\(T_{c'}z-T_cz\) as the boundary \(dK_mz\).
\end{proof}

Let \(V\) now be a finite-dimensional vector space
over \(k\).  Write \(\mathbb P(V)\) for its set
of one-dimensional subspaces; it is empty for \(V=0\).
Choose one nonzero vector in each line.  This defines
a homogeneous function \(\lambda:V\setminus\{0\}\to k^\times\)
by expressing a vector as \(\lambda(v)\) times the
chosen vector in its line.  Thus
\(\lambda(av)=a\lambda(v)\) for \(a\ne0\).

For a representation \(\rho:G\to\operatorname{GL}(V)\),
use the right action \(\ell g=\rho(g)^{-1}\ell\),
and put
\[
                  c_\rho(g,\ell)
                    =\frac{\lambda(v)}{\lambda(\rho(g)^{-1}v)},
                    \qquad 0\ne v\in\ell.
\]
Homogeneity makes this independent of \(v\).
It satisfies the cocycle identity in
Lemma~\ref{sf-transfer}.
Changing \(\lambda\) gives precisely the coboundary
in that lemma, with
\(b(\ell)=\lambda(v)/\lambda'(v)\).

\begin{proposition}\label{sf-additive-transfer}
Let \(A\) be a commutative ring in which \(Q=1\),
and let \(t:C_m(k^\times;A)\to A\) be a cocycle
of positive degree.  For every finite-dimensional
\(k\)-space \(V\), define
\[
                    \theta_V=t\circ T_{c_{\mathrm{id}}}.
\]
These functionals vanish on boundaries.  Their values
on cycles are invariant under linear isomorphisms.
For representations \(\rho\) on \(V\) and \(\sigma\)
on \(W\), and any cycle \(z\in C_m(G;A)\),
\[
 \theta_{V\oplus W}((\rho\oplus\sigma)_*z)
       =\theta_V(\rho_*z)+\theta_W(\sigma_*z).
\]
In dimension one, the scalar representation identifies
\(\theta_k\) with \(t\) on cycles.
\end{proposition}
\begin{proof}
The chain-map property proves the assertion about
boundaries.  Pulling the homogeneous function through
a linear isomorphism gives identical summands under
the resulting bijection of projective lines.
Any other choice gives the same value on cycles by
Lemma~\ref{sf-transfer}.

For the direct sum, choose
\[
 \lambda(v,w)=
 \begin{cases}
     \lambda_V(v),&v\ne0,\\
     \lambda_W(w),&v=0.
 \end{cases}
\]
Send a line represented by \((v,w)\) to \(kv\in\mathbb P(V)\)
when \(v\ne0\), and otherwise to \(kw\in\mathbb P(W)\).
This map is equivariant for \(\rho\oplus\sigma\).
The scalar cocycle is the pullback of the one on
the appropriate summand.
A fiber over \(kv\in\mathbb P(V)\), after fixing
the nonzero representative \(v\), consists uniquely
of the lines \(k(v,w)\) with \(w\in W\).
It has \(Q^{\dim W}\) elements.
A fiber over a line in \(W\) has one element.
Consequently the transfer for the direct sum is
\(Q^{\dim W}T_\rho+T_\sigma=T_\rho+T_\sigma\)
with coefficients in \(A\).
Choice independence proves the claimed formula for
the originally chosen functionals on cycles.

For \(V=k\), there is one projective line.
Choosing \(\lambda(v)=v\) gives
\(c(g)=g\), so its transfer is the identity.
\end{proof}

\subsection{The determinant value}

\begin{lemma}\label{sf-determinant}
For \(g\in\operatorname{GL}(V)\),
\[
                 \prod_{\ell\in\mathbb P(V)}c(g,\ell)
                              =\det(g)\quad\text{in }k^\times.
\]
\end{lemma}
\begin{proof}
Denote the product by \(D(g)\).
The cocycle identity and the permutation of the
projective lines imply \(D(gh)=D(g)D(h)\).
Changing \(\lambda\) does not change the product,
because all factors \(b(\ell)^{-1}b(\ell g)\) cancel.

Choose coordinates and take \(\lambda(v)\) to be
the first nonzero coordinate of \(v\).
For \(g=\operatorname{diag}(a_1,\ldots,a_n)\), a line
whose first nonzero coordinate has index \(i\)
contributes \(a_i\).  There are \(Q^{n-i}\) such
lines, obtained by making that coordinate equal
to \(1\).  Since \(a_i^Q=a_i\) in the finite field,
\[
                      D(g)=\prod_i a_i^{Q^{n-i}}
                            =\prod_i a_i=\det(g).
\]
If \(Q=2\), the target group \(k^\times\) is
trivial and the result already follows.
Otherwise choose \(a\in k\setminus\{0,1\}\).
For \(i\ne j\), let \(T_{ij}(s)=I+sE_{ij}\),
and let \(d\) be diagonal with entry \(a\) at
\(i\) and \(1\) elsewhere.  Then
\[
      dT_{ij}(s)d^{-1}T_{ij}(s)^{-1}=T_{ij}((a-1)s).
\]
Every transvection is therefore a commutator, so
the homomorphism \(D\) is \(1\) on transvections.
Gaussian elimination writes every invertible matrix
as a product of diagonal matrices and transvections:
if a pivot is zero, add a row having a nonzero entry
in that column, then clear the remaining column
entries and continue.  Row scalings supply the
diagonal matrices.
Both \(D\) and determinant agree on these generators,
which proves the lemma.  In dimension zero both
products are \(1\).
\end{proof}

\begin{theorem}\label{sf-functional}
Under the hypotheses of Proposition~\ref{sf-cyclic-functional},
there are functionals
\[
             \theta_V:C_{2E+1}(\operatorname{GL}(V);R)\to R
\]
which vanish on boundaries, are invariant under
linear isomorphisms on cycles, and are additive
under direct sums on cycles.  They satisfy
\[
 \theta_V\bigl(P_g(\chi_{x,V})_*z_E\bigr)
                         =x^E v(\det g),
 \qquad
 \chi_{x,V}(j)=\zeta^{xj}I_V.
\]
\end{theorem}
\begin{proof}
Since \(q\mid Q-1\), we have \(Q=1\) in \(R\).
Apply Proposition~\ref{sf-additive-transfer} to the
cyclic cocycle \(t\).
Scalar matrices act trivially on \(\mathbb P(V)\),
and their scalar cocycle is the scalar itself.
Thus, directly from the tuple formula for transfer,
\[
 T\bigl(P_g(\chi_{x,V})_*z_E\bigr)
     =\sum_{\ell\in\mathbb P(V)}
          P_{c(g,\ell)}(\chi_x)_*z_E.
\]
Indeed, before the inserted \(g\), all projective
lines remain fixed by the preceding scalar entries;
after it, the line may change, but the remaining
scalar cocycle entries do not depend on that line.
Apply Proposition~\ref{sf-cyclic-functional} and then
Lemma~\ref{sf-determinant} to obtain
\[
 x^E\sum_\ell v(c(g,\ell))
       =x^E v\!\left(\prod_\ell c(g,\ell)\right)
       =x^E v(\det g).
\]
If the coefficient modulus is \(q=1\), the same
assertion holds with all functionals zero; no
nontrivial root of unity is then required.
\end{proof}

Taking \(E=e+1\) proves the source assertions
\(\mathrm{SCoc}\), \(\mathrm{STrans}\), and
\(\mathrm{STheta}\), with their stated coefficients
and normalization.  The determinant calculation also
explains why no exceptional claim about
\(\operatorname{GL}_2(\mathbb F_2)\) is needed:
its multiplicative residue group is trivial.

\begin{corollary}\label{sf-primitive-family}
Let \(A\) be a commutative coefficient ring, let
\(\theta\) be any positive-degree additive family
with values in \(A\) as above, let \(B\) be a commutative ring, and let
\(r:B\to k\) be a ring homomorphism.
Composing reduction of matrices with \(\theta\)
gives a family on \(\operatorname{GL}_n(B)\) which
vanishes on boundaries, is invariant on cycles under
\(g\mapsto\operatorname{diag}(g,I)\), and is
multiplied by \(d\) under
\(g\mapsto\operatorname{diag}(g,\ldots,g)\)
with \(d\) copies.
\end{corollary}
\begin{proof}
Reduction is a group homomorphism and therefore
commutes with the bar boundary.
The contribution of a trivial representation on
a positive-degree cycle is zero.
To see this without an acyclicity assumption,
the bar chain group of the trivial group is \(A\)
in every degree; its boundary is the identity in
positive even degrees and zero in odd degrees.
Its positive-degree homology is consequently zero.
The image of the cycle under a trivial representation
is therefore a boundary, annihilated by \(\theta\).
Direct-sum additivity now proves stabilization
invariance, and iteration proves multiplication
by \(d\).  For \(d=0\) the representation is trivial
on the zero space and gives zero.
\end{proof}

This corollary is exactly the source assertion
\(\mathrm{SPhi}\).  The finite-field evaluations and
their stable additivity are now proved.  Constructing
the integral arithmetic cycles with the prescribed
evaluations, and bounding their possible values by
the stable rational homology calculation, remain
separate tasks.


